\chapter{Planificación y análisis de costes}
\label{chap:planificacionycostes}

\lettrine{E}{n} este capítulo se presenta la planificación aplicada durante el desarrollo del proyecto y un análisis de los costes asociados. Se describe la planificación inicial, las sprints que han estructurado el trabajo y las desviaciones observadas en el seguimiento.

\section{Planificación}
\label{sec:planificacion}

La planificación inicial se definió en 9 sprints, cada una organizada internamente siguiendo el modelo clásico de cuatro fases: análisis, diseño, implementación y pruebas. La estimación de esfuerzo inicial se basó en jornadas de 6 horas diarias y en una duración prevista del proyecto de cuatro meses. No obstante, las complejidades técnicas y errores imprevistos motivaron desviaciones en varias iteraciones, tal y como se comenta en la sección de seguimiento.

\begin{table}[H]
  \centering
  \resizebox{\textwidth}{!}{
    \rowcolors{2}{white}{udcgray!25}
    \begin{tabular}{l|c|c|c}
      \rowcolor{udcpink!25}
      \textbf{Tarea} & \textbf{Fecha inicio} & \textbf{Días estimados} & \textbf{Días reales} \\

      Definición tipos Emergency & 04/03 & 0.5 & 0.5 \\
      Definición modelo Assignment & 04/03 & 0.5 & 0.5 \\
      Definición contratos OpenApi & 04/03 & 1 & 1 \\
      Diseño entidades JPA & 05/03 & 0.5 & 0.5 \\

      Implementación entidad Emergency & 06/03 & 0.5 & 0.5 \\
      Implementación entidad Assignment & 06/03 & 0.5 & 0.5 \\
      Implementación repositorios & 06/03 & 0.5 & 0.5 \\

      UseCase Emergency & 08/03 & 0.5 & 0.5 \\
      UseCase Assignment & 08/03 & 0.5 & 0.5 \\

      Endpoints /emergencies & 09/03 & 0.5 & 0.5 \\
      Endpoints /assignments & 09/03 & 0.5 & 0.5 \\

      Tests unitarios backend & 11/03 & 1 & 1 \\
      Tests integración backend & 11/03 & 1 & 1 \\

      Consumir Api desde frontend & 10/03 & 0.5 & 0.5 \\
      Vista mapa (crear Emergency) & 11/03 & 1 & 1 \\
      Formulario Emergency & 11/03 & 0.5 & 0.5 \\
      Listado Emergencies & 12/03 & 0.5 & 0.5 \\
      Actualización UI Assignment & 12/03 & 1 & 1 \\
      Listado Assignments & 13/03 & 0.5 & 0.5 \\

      Consumir Api desde Android & 07/03 & 0.5 & 0.5 \\
      Lista Assignments Android & 12/03 & 1 & 1 \\
      Detalle Assignment Android & 13/03 & 0.5 & 0.5 \\
      Corrección bugs & 15/03 & 1 & 1 \\

    \end{tabular}
  }
  \caption{Planificación del Sprint 3: Emergencias y asignaciones básicas}
  \label{tab:sprint3}
\end{table}

En la tabla~\ref{tab:sprint3} se puede observar un desglose del sprint 3, que se centró en la implementación de la gestión de emergencias y asignaciones básicas. Cada tarea se estimó inicialmente en función de su complejidad, esto sirvió para poder observar desviaciones y ajustar la planificación de los siguientes sprints. En este caso, el sprint se completó dentro del plazo previsto, lo que permitió mantener el ritmo de desarrollo planificado.


\section{Sprints}
\label{sec:sprints}

El proyecto se estructuró en una sucesión de sprints, cada uno de los cuales incrementaba la funcionalidad de los artefactos entregados.  A continuación se detalla cada sprint con sus objetivos y resultados principales.

\subsection{Sprint 0: Contextualización}
En esta fase inicial se realizó la contextualización del dominio y la justificación del valor del \acrshort{tfm}. Se investigaron fuentes, alternativas a la aplicación y tecnologías potenciales (mapas, sistema de notificaciones android, sistema de recomendación y reglas básico). El objetivo fue delimitar el alcance y dominio del proyecto para que este sea viable en un espacio temporal acotado.

\subsection{Sprint 1: Esqueleto y arquitectura base de aplicación móvil}
Se creó el esqueleto del proyecto móvil y se definieron las dependencias y la arquitectura base (estructura de paquetes, patrones de navegación, y la configuración inicial de Jetpack Compose). Además se implementaron los módulos iniciales de autenticación y sincronización con el Backend (endpoints básicos) para poder iterar sobre funcionalidades reales.

\subsection{Sprint 2: Actualizar la organización de recursos}
Se introdujo el modelo unificado de recursos y organizaciones preservando la compatibilidad con las entidades Legacy (Equipo y Vehículo). Esto incluyó definir atributos de capacidad, áreas operativas (cuadrantes) y las relaciones necesarias a nivel de base de datos y DTOs. Se arreglaron los endpoints antiguos y DTO's parara la gestión CRUD de organizaciones y sus recursos.

\subsection{Sprint 3: Emergencias y asignaciones básicas}
Se implementó la gestión de emergencias con soporte para dos tipos de localización: punto y cuadrante. Se añadió la creación y edición de Emergencias, la definición de cuadrantes y una primera versión del flujo de creación de Asignaciones. Esta iteración incluyó validaciones de consistencia (ej. comprobar que un cuadrante pertenece a una emergencia) y pruebas de integración explicitas para el servicio de Asignaciones.

\subsection{Sprint 4: Nuevo sistema de logs operativos}

Se definió un nuevo esquema para los logs operativos, ahora centralizados en las Asignaciones de recursos en las emergencias. El trabajo abarcó el nuevo modelado y la persistencia que permiten conservar el histórico.

\subsection{Sprint 5: Notificaciones en Android y aceptación de asignaciones}
Se integró Firebase Cloud Messaging (\acrshort{fcm}) en la aplicación Android: registro de MobileDevice, manejo de tokens y recepción de notificaciones. Se implementó el flujo de envío de notificaciones desde el Backend y la experiencia de usuario en la app para aceptar o rechazar asignaciones recibidas.

\subsection{Sprint 6: Protocolos semiautomáticos (motor de reglas simple)}
Se introdujo la capacidad de definir Protocols almacenados en \acrshort{jsonb} y un motor de reglas ligero que permite generar propuestas de Asignaciones de forma semiautomática en el frontend. Esto incluyó la definición del formato JSON para los protocolos, la implementación de un evaluador simple y la integración con el pipeline de creación de propuestas.

\subsection{Sprint 7 y 8: Finalización, pruebas y despliegue}
Las iteraciones finales se dedicaron a la estabilización: pruebas end-to-end, ajuste de UX tanto del frontal como de la aplicación Android, corrección de errores, poblamiento de datos de ejemplo y preparación de artefactos de despliegue (contenedores Docker para BBDD y servicio). También incluyeron la redacción y revisión de la memoria y la generación de la documentación.

\section{Seguimiento}
\label{sec:seguimiento}

El seguimiento se llevó a cabo mediante reuniones mensuales con el tutor y el uso de un tablero Kanban de tareas que permitía visualizar el estado de ellas. En cada revisión se presentaron artefactos funcionales (pantallas, endpoints, wireframes).

Las desviaciones más relevantes provinieron de problemas técnicos con la representación de cuadrantes en el mapa, configuración de la aplicación Android y el motor de reglas básico para los protocolos. 

\section{Análisis de costes}
\label{sec:costes}

El análisis de costes se ha planteado distinguiendo entre costes directos de personal, costes indirectos de infraestructura y un margen de contingencia para imprevistos. Para la valoración económica del esfuerzo se han tomado como referencia salarios publicados por Talent.com para España, donde un perfil de \textit{Programador} se sitúa en 27.500~\euro/año (14,10~\euro/h)  \cite{salario_programador}y un perfil de \textit{Analista programador} en 29.120~\euro/año (14,93~\euro/h)~\cite{salario_analista_programador}. Estos valores resultan razonables para un TFM técnico con tareas de análisis, desarrollo, pruebas e integración.

\subsection{Estimación del esfuerzo}

La planificación se ha estructurado en 7 sprints. Si se asume una duración media de 2 semanas por sprint, el proyecto cubre aproximadamente 14 semanas de trabajo efectivo. Considerando una dedicación máxima de 6 horas al día y trabajo únicamente de lunes a viernes, el esfuerzo total estimado se calcula como sigue:

\[
7\ \text{sprints} \times 2\ \text{semanas/sprint} \times 5\ \text{días/semana} \times 6\ \text{h/día} = 420\ \text{horas}
\]

Este valor representa una estimación prudente del esfuerzo total del proyecto. La distribución del tiempo no es homogénea entre sprints, ya que las primeras iteraciones requieren más análisis y diseño, mientras que las finales concentran más trabajo de integración, pruebas, corrección de errores y documentación.

\subsection{Costes directos de personal}

Para la valoración del trabajo se ha tomado como base el perfil de \textit{Analista programador}, al ser el que mejor refleja un proyecto de estas características, donde una misma persona asume tareas de definición funcional, implementación y validación. Con la tarifa media indicada por la fuente utilizada, el coste directo estimado es:

\[
420\ \text{h} \times 14,93\ \text{\euro/h} = 6.270,60\ \text{\euro}
\]

De forma complementaria, si se valorase el trabajo con el perfil de \textit{Programador}, el coste sería:

\[
420\ \text{h} \times 14,10\ \text{\euro/h} = 5.922,00\ \text{\euro}
\]

Para reflejar ambas referencias, se adopta como coste base del proyecto el punto medio entre ambos perfiles, es decir, \textbf{6.096,30~\euro}, que representa mejor la mezcla de tareas de programación y análisis realizada durante el desarrollo.

\subsection{Costes indirectos}

A los costes de personal se añaden los costes indirectos de ejecución. Para este tipo de proyecto, los más relevantes son la amortización del equipo de trabajo y los suministros básicos asociados al desarrollo:

\begin{itemize}
  \item Amortización de hardware: se considera el uso de un equipo portátil de desarrollo valorado en 900~\euro, con una vida útil de 4 años. Para un proyecto de unos 3,5 meses, el coste imputable es de aproximadamente 66,00~\euro.
  \item Conectividad y suministros: se estiman 40~\euro/mes de internet y 30~\euro/mes de electricidad proporcional al tiempo de desarrollo. Para 3,5 meses, esto supone 245,00~\euro.
\end{itemize}

\subsection{Resumen económico}

\begin{table}[H]
  \centering
  \rowcolors{2}{white}{udcgray!25}
  \resizebox{1\textwidth}{!}{%
    \begin{tabular}{c|c|c|c|c}
      \rowcolor{udcpink!25}
    Concepto & Importe estimado (\euro) & Observaciones \\
    Coste directo de personal & 6.096,30 & Media entre programador y analista programador \\
    Amortización de hardware & 66,00 & Portátil de desarrollo imputado al proyecto \\
    Conectividad y suministros & 245,00 & Internet y electricidad proporcionales \\
    \hline
    \textbf{Total estimado} & \textbf{6.407,3} & \\
    \hline
   \end{tabular}%
  }
  \caption{Desglose de costes estimados del proyecto}
  \label{tab:costes_resumen}
\end{table}

La tabla~\ref{tab:costes_resumen} muestra que el mayor peso del presupuesto corresponde al trabajo humano, seguido por los costes de infraestructura y suministros. Esto es coherente con un proyecto de software de ámbito académico, donde la inversión principal no está en equipamiento especializado sino en la dedicación de horas de desarrollo, integración y validación.


\subsection{Conclusión del análisis de costes}

El coste final del proyecto asciende a \textbf{6.407,3\euro} con una dedicación total de \textbf{420 horas}. Estas cifras reflejan el alcance y la coste del \acrshort{tfm}, considerando tanto el esfuerzo humano como los recursos materiales necesarios para su ejecución. Es importante destacar que esta estimación se ha realizado con criterios prudentes y basándose en datos salariales reales, lo que aporta una visión más ajustada del valor económico del trabajo realizado en este proyecto.
